\SourceLeaf{389}%
to 1 to 40000, and for the engraving to 1 to 80000. On this basis, the number of sheets (of the same size as above) was estimated at 208, the time required at 15 years, the cost of the work, the engraving and the printing of 3000 copies at 4232000 francs, and finally the selling price of each sheet at 7 francs 50 centimes, or 1560 francs for the whole map. The entire work falls within the remit of the Depot, but it is counting on the cooperation of the cadastral service, though the report does not make it entirely clear to what extent. It seems that all the detailed survey work is expected from this authority (which is responsible to another ministry), so that only the trigonometric work and the determination of elevations fall to the engineer-geographers in the survey. This dependence on another authority, with which cooperation seems less than harmonious (the failure of the first plan is attributed to insufficient support from the cadastral service), might prove very detrimental to the rapid progress hoped for in this great undertaking. Once the work is complete, a large compendium is also to be issued, making public both the numerical results for the positions and elevations of the trigonometric points and all the measurements on which they are based. This would of course satisfy all our wishes. But there is reason to fear that the time may still be very far off, and the first-order triangles and triangulation points in particular are of the greatest interest even for more advanced scientific purposes. We therefore cannot suppress the earnest wish that the full publication of the first-order triangles (all of which have already been measured) should be delayed no longer and that these results should be made available at once for the benefit of science. Those interested in higher geodesy would be all the more grateful if future volumes of the Memorial fulfilled these wishes. since the volumes published so far, which are so substantial in other respects, have given them the least attention, and the few theoretical dissertations and auxiliary tables offer no compensation for the lack of \textit{facts} that the Depot could supply so abundantly.
\begin{center}\rule{.3\textwidth}{.3pt}\end{center}
